% ------------------------------------------------------------------------
%  Energy stage cost: CPU vs naive GPU sum vs FMM.
%  Generated by scripts/plot_energy_methods.py -- regenerate rather than edit.
%  Include with % ------------------------------------------------------------------------
%  Energy stage cost: CPU vs naive GPU sum vs FMM.
%  Generated by scripts/plot_energy_methods.py -- regenerate rather than edit.
%  Include with % ------------------------------------------------------------------------
%  Energy stage cost: CPU vs naive GPU sum vs FMM.
%  Generated by scripts/plot_energy_methods.py -- regenerate rather than edit.
%  Include with % ------------------------------------------------------------------------
%  Energy stage cost: CPU vs naive GPU sum vs FMM.
%  Generated by scripts/plot_energy_methods.py -- regenerate rather than edit.
%  Include with \input{figures/energy_methods_cluster}
%  Needs \usepackage{pgfplots} and \usetikzlibrary{arrows.meta} in the preamble.
% ------------------------------------------------------------------------
\input{figures/data/energy_methods_cluster_ref}

\pgfplotsset{compat=1.16}

\definecolor{energycpu}{HTML}{D95F02}
\definecolor{energynaive}{HTML}{9DC3E0}
\definecolor{energyfmm}{HTML}{2C5F8D}

\begin{figure}[!htb]
\centering
\begin{tikzpicture}
\begin{axis}[
    ybar,
    % 15cm against the 16.03cm textwidth of the 12pt/2.5cm page. Measured, not
    % assumed: 15.4 still fits, 13.6 gains nothing.
    width=15.0cm, height=7.4cm,
    bar width=13pt,
    enlarge x limits=0.30,
    % A shared linear axis would put every bar but the slowest on the baseline:
    % the times span 5756x. log origin y=infty is what makes the bars start at
    % ymin rather than at 1, which on a log axis is where ybar puts them by
    % default -- without it the sub-second bars would hang below the frame.
    ymode=log, log origin y=infty,
    ymin=0.05, ymax=20000,
    ytick={0.1,1,10,100,1000,10000},
    % 1, 10, 100 rather than 10^0, 10^1, 10^2. The times are wall clock and get
    % read off as such; the exponents make a reader translate before comparing.
    log ticks with fixed point,
    ylabel={energy stage time [s]},
    ylabel style={font=\footnotesize},
    symbolic x coords={1VSZ,6VYB,H1N1},
    xtick=data,
    % The atom count belongs on the axis, not in the caption: the three columns
    % are three problem sizes, and without it the H1N1 column reads as the same
    % work taking longer rather than as two orders of magnitude more of it.
    xticklabels={{1VSZ\\{\scriptsize 14\,299 atoms}},{6VYB\\{\scriptsize 44\,495 atoms}},{H1N1\\{\scriptsize 14\,442\,610 atoms}}},
    xticklabel style={font=\footnotesize, align=center},
    tick label style={font=\scriptsize},
    ymajorgrids=true,
    major grid style={black!12},
    % Clipping stays ON. log origin y=infty gives every bar a base far below the
    % frame -- a large finite log coordinate, not a real infinity -- and the clip
    % is the only thing keeping those bases off the page. The DNF annotation
    % therefore cannot be drawn in here; it is anchored to a \coordinate and
    % drawn after \end{axis}, where clipping no longer applies.
    % -0.22, not the -0.14 the other figures use: the tick labels here are two
    % lines, and the legend has to clear the atom count as well as the name.
    legend style={
      at={(0.5,-0.22)}, anchor=north, draw=none,
      legend columns=3, font=\footnotesize, column sep=0.8em,
    },
    error bars/error bar style={black!85, line width=0.8pt},
    error bars/error mark options={
      mark size=3pt, rotate=90, line width=0.8pt, black!85,
    },
    % The white backing is load-bearing on the H1N1 group: a four-digit label is
    % wider than a 13pt bar, so it overhangs onto the neighbouring bar, which
    % there is a full-height block of colour.
    every node near coord/.append style={
      font=\tiny, color=black!55, anchor=south, yshift=4pt,
      inner sep=1pt, fill=white, fill opacity=0.85, text opacity=1,
    },
]
  \addplot+[
    ybar, bar shift=-14.0pt, fill=energycpu, draw=energycpu!70!black, line width=0.4pt,
    nodes near coords, point meta=explicit symbolic,
    error bars/.cd, y dir=both, y explicit,
  ] table[x=molecule, y=cpu, y error minus=cpu_em, y error plus=cpu_ep,
          meta=cpu_lbl] {figures/data/energy_methods_cluster.dat};
  \addlegendentry{CPU}
  \addplot+[
    ybar, bar shift=0.0pt, fill=energynaive, draw=energynaive!70!black, line width=0.4pt,
    nodes near coords, point meta=explicit symbolic,
    error bars/.cd, y dir=both, y explicit,
  ] table[x=molecule, y=naive, y error minus=naive_em, y error plus=naive_ep,
          meta=naive_lbl] {figures/data/energy_methods_cluster.dat};
  \addlegendentry{naive GPU $O(NM)$}
  \addplot+[
    ybar, bar shift=14.0pt, fill=energyfmm, draw=energyfmm!70!black, line width=0.4pt,
    nodes near coords, point meta=explicit symbolic,
    error bars/.cd, y dir=both, y explicit,
  ] table[x=molecule, y=fmm, y error minus=fmm_em, y error plus=fmm_ep,
          meta=fmm_lbl] {figures/data/energy_methods_cluster.dat};
  \addlegendentry{FMM}

  \coordinate (dnf-H1N1-cpu) at ([xshift=-14.0pt]axis cs:H1N1,20000);

\end{axis}
\draw[-{Stealth[length=2.6mm,width=2.2mm]}, line width=1pt, energycpu!75!black]
  (dnf-H1N1-cpu) -- ++(0,0.62cm)
  node[above, inner sep=1pt, font=\scriptsize\bfseries, text=black!70]{DNF};

\end{tikzpicture}

\caption[Energy stage cost by method]{\textbf{Cost of the energy stage on
    \energyMoleculesenergymethodscluster{}.} Median wall time of the energy calculation stage on a
    logarithmic axis; the number above each bar is that median in
    seconds.\energyMachineenergymethodscluster\energyCaveatsenergymethodscluster}
\label{fig:energy-methods-cluster}
\end{figure}

%  Needs \usepackage{pgfplots} and \usetikzlibrary{arrows.meta} in the preamble.
% ------------------------------------------------------------------------
% generated by scripts/plot_energy_methods.py -- do not edit
\def\energyMoleculesenergymethodscluster{1VSZ, 6VYB, H1N1}
\def\energyRepeatsenergymethodscluster{3}
\def\energyMachineenergymethodscluster{ Measured on an A100 node of the TU Berlin HPC (Table~\ref{tab:test-systems}) with 16 MPI ranks.}
\def\energyFmmConfigenergymethodscluster{ The FMM ran at $P = 9$, $\theta = 0.3$, $n_{leaf,src} = 16$, $n_{leaf,tgt} = 1024$ for 1VSZ and 6VYB; $P = 10$, $\theta = 0.4$, $n_{leaf} = 256$ for H1N1.}
\def\energyReferenceenergymethodscluster{the CPU path where it ran and the naive GPU sum otherwise}
\def\energyCaveatsenergymethodscluster{ H1N1: measured on four H200 GPUs of the TU Berlin HPC (Table~\ref{tab:test-systems}) with 8 MPI ranks; 1 run per method, not 3.}


\pgfplotsset{compat=1.16}

\definecolor{energycpu}{HTML}{D95F02}
\definecolor{energynaive}{HTML}{9DC3E0}
\definecolor{energyfmm}{HTML}{2C5F8D}

\begin{figure}[!htb]
\centering
\begin{tikzpicture}
\begin{axis}[
    ybar,
    % 15cm against the 16.03cm textwidth of the 12pt/2.5cm page. Measured, not
    % assumed: 15.4 still fits, 13.6 gains nothing.
    width=15.0cm, height=7.4cm,
    bar width=13pt,
    enlarge x limits=0.30,
    % A shared linear axis would put every bar but the slowest on the baseline:
    % the times span 5756x. log origin y=infty is what makes the bars start at
    % ymin rather than at 1, which on a log axis is where ybar puts them by
    % default -- without it the sub-second bars would hang below the frame.
    ymode=log, log origin y=infty,
    ymin=0.05, ymax=20000,
    ytick={0.1,1,10,100,1000,10000},
    % 1, 10, 100 rather than 10^0, 10^1, 10^2. The times are wall clock and get
    % read off as such; the exponents make a reader translate before comparing.
    log ticks with fixed point,
    ylabel={energy stage time [s]},
    ylabel style={font=\footnotesize},
    symbolic x coords={1VSZ,6VYB,H1N1},
    xtick=data,
    % The atom count belongs on the axis, not in the caption: the three columns
    % are three problem sizes, and without it the H1N1 column reads as the same
    % work taking longer rather than as two orders of magnitude more of it.
    xticklabels={{1VSZ\\{\scriptsize 14\,299 atoms}},{6VYB\\{\scriptsize 44\,495 atoms}},{H1N1\\{\scriptsize 14\,442\,610 atoms}}},
    xticklabel style={font=\footnotesize, align=center},
    tick label style={font=\scriptsize},
    ymajorgrids=true,
    major grid style={black!12},
    % Clipping stays ON. log origin y=infty gives every bar a base far below the
    % frame -- a large finite log coordinate, not a real infinity -- and the clip
    % is the only thing keeping those bases off the page. The DNF annotation
    % therefore cannot be drawn in here; it is anchored to a \coordinate and
    % drawn after \end{axis}, where clipping no longer applies.
    % -0.22, not the -0.14 the other figures use: the tick labels here are two
    % lines, and the legend has to clear the atom count as well as the name.
    legend style={
      at={(0.5,-0.22)}, anchor=north, draw=none,
      legend columns=3, font=\footnotesize, column sep=0.8em,
    },
    error bars/error bar style={black!85, line width=0.8pt},
    error bars/error mark options={
      mark size=3pt, rotate=90, line width=0.8pt, black!85,
    },
    % The white backing is load-bearing on the H1N1 group: a four-digit label is
    % wider than a 13pt bar, so it overhangs onto the neighbouring bar, which
    % there is a full-height block of colour.
    every node near coord/.append style={
      font=\tiny, color=black!55, anchor=south, yshift=4pt,
      inner sep=1pt, fill=white, fill opacity=0.85, text opacity=1,
    },
]
  \addplot+[
    ybar, bar shift=-14.0pt, fill=energycpu, draw=energycpu!70!black, line width=0.4pt,
    nodes near coords, point meta=explicit symbolic,
    error bars/.cd, y dir=both, y explicit,
  ] table[x=molecule, y=cpu, y error minus=cpu_em, y error plus=cpu_ep,
          meta=cpu_lbl] {figures/data/energy_methods_cluster.dat};
  \addlegendentry{CPU}
  \addplot+[
    ybar, bar shift=0.0pt, fill=energynaive, draw=energynaive!70!black, line width=0.4pt,
    nodes near coords, point meta=explicit symbolic,
    error bars/.cd, y dir=both, y explicit,
  ] table[x=molecule, y=naive, y error minus=naive_em, y error plus=naive_ep,
          meta=naive_lbl] {figures/data/energy_methods_cluster.dat};
  \addlegendentry{naive GPU $O(NM)$}
  \addplot+[
    ybar, bar shift=14.0pt, fill=energyfmm, draw=energyfmm!70!black, line width=0.4pt,
    nodes near coords, point meta=explicit symbolic,
    error bars/.cd, y dir=both, y explicit,
  ] table[x=molecule, y=fmm, y error minus=fmm_em, y error plus=fmm_ep,
          meta=fmm_lbl] {figures/data/energy_methods_cluster.dat};
  \addlegendentry{FMM}

  \coordinate (dnf-H1N1-cpu) at ([xshift=-14.0pt]axis cs:H1N1,20000);

\end{axis}
\draw[-{Stealth[length=2.6mm,width=2.2mm]}, line width=1pt, energycpu!75!black]
  (dnf-H1N1-cpu) -- ++(0,0.62cm)
  node[above, inner sep=1pt, font=\scriptsize\bfseries, text=black!70]{DNF};

\end{tikzpicture}

\caption[Energy stage cost by method]{\textbf{Cost of the energy stage on
    \energyMoleculesenergymethodscluster{}.} Median wall time of the energy calculation stage on a
    logarithmic axis; the number above each bar is that median in
    seconds.\energyMachineenergymethodscluster\energyCaveatsenergymethodscluster}
\label{fig:energy-methods-cluster}
\end{figure}

%  Needs \usepackage{pgfplots} and \usetikzlibrary{arrows.meta} in the preamble.
% ------------------------------------------------------------------------
% generated by scripts/plot_energy_methods.py -- do not edit
\def\energyMoleculesenergymethodscluster{1VSZ, 6VYB, H1N1}
\def\energyRepeatsenergymethodscluster{3}
\def\energyMachineenergymethodscluster{ Measured on an A100 node of the TU Berlin HPC (Table~\ref{tab:test-systems}) with 16 MPI ranks.}
\def\energyFmmConfigenergymethodscluster{ The FMM ran at $P = 9$, $\theta = 0.3$, $n_{leaf,src} = 16$, $n_{leaf,tgt} = 1024$ for 1VSZ and 6VYB; $P = 10$, $\theta = 0.4$, $n_{leaf} = 256$ for H1N1.}
\def\energyReferenceenergymethodscluster{the CPU path where it ran and the naive GPU sum otherwise}
\def\energyCaveatsenergymethodscluster{ H1N1: measured on four H200 GPUs of the TU Berlin HPC (Table~\ref{tab:test-systems}) with 8 MPI ranks; 1 run per method, not 3.}


\pgfplotsset{compat=1.16}

\definecolor{energycpu}{HTML}{D95F02}
\definecolor{energynaive}{HTML}{9DC3E0}
\definecolor{energyfmm}{HTML}{2C5F8D}

\begin{figure}[!htb]
\centering
\begin{tikzpicture}
\begin{axis}[
    ybar,
    % 15cm against the 16.03cm textwidth of the 12pt/2.5cm page. Measured, not
    % assumed: 15.4 still fits, 13.6 gains nothing.
    width=15.0cm, height=7.4cm,
    bar width=13pt,
    enlarge x limits=0.30,
    % A shared linear axis would put every bar but the slowest on the baseline:
    % the times span 5756x. log origin y=infty is what makes the bars start at
    % ymin rather than at 1, which on a log axis is where ybar puts them by
    % default -- without it the sub-second bars would hang below the frame.
    ymode=log, log origin y=infty,
    ymin=0.05, ymax=20000,
    ytick={0.1,1,10,100,1000,10000},
    % 1, 10, 100 rather than 10^0, 10^1, 10^2. The times are wall clock and get
    % read off as such; the exponents make a reader translate before comparing.
    log ticks with fixed point,
    ylabel={energy stage time [s]},
    ylabel style={font=\footnotesize},
    symbolic x coords={1VSZ,6VYB,H1N1},
    xtick=data,
    % The atom count belongs on the axis, not in the caption: the three columns
    % are three problem sizes, and without it the H1N1 column reads as the same
    % work taking longer rather than as two orders of magnitude more of it.
    xticklabels={{1VSZ\\{\scriptsize 14\,299 atoms}},{6VYB\\{\scriptsize 44\,495 atoms}},{H1N1\\{\scriptsize 14\,442\,610 atoms}}},
    xticklabel style={font=\footnotesize, align=center},
    tick label style={font=\scriptsize},
    ymajorgrids=true,
    major grid style={black!12},
    % Clipping stays ON. log origin y=infty gives every bar a base far below the
    % frame -- a large finite log coordinate, not a real infinity -- and the clip
    % is the only thing keeping those bases off the page. The DNF annotation
    % therefore cannot be drawn in here; it is anchored to a \coordinate and
    % drawn after \end{axis}, where clipping no longer applies.
    % -0.22, not the -0.14 the other figures use: the tick labels here are two
    % lines, and the legend has to clear the atom count as well as the name.
    legend style={
      at={(0.5,-0.22)}, anchor=north, draw=none,
      legend columns=3, font=\footnotesize, column sep=0.8em,
    },
    error bars/error bar style={black!85, line width=0.8pt},
    error bars/error mark options={
      mark size=3pt, rotate=90, line width=0.8pt, black!85,
    },
    % The white backing is load-bearing on the H1N1 group: a four-digit label is
    % wider than a 13pt bar, so it overhangs onto the neighbouring bar, which
    % there is a full-height block of colour.
    every node near coord/.append style={
      font=\tiny, color=black!55, anchor=south, yshift=4pt,
      inner sep=1pt, fill=white, fill opacity=0.85, text opacity=1,
    },
]
  \addplot+[
    ybar, bar shift=-14.0pt, fill=energycpu, draw=energycpu!70!black, line width=0.4pt,
    nodes near coords, point meta=explicit symbolic,
    error bars/.cd, y dir=both, y explicit,
  ] table[x=molecule, y=cpu, y error minus=cpu_em, y error plus=cpu_ep,
          meta=cpu_lbl] {figures/data/energy_methods_cluster.dat};
  \addlegendentry{CPU}
  \addplot+[
    ybar, bar shift=0.0pt, fill=energynaive, draw=energynaive!70!black, line width=0.4pt,
    nodes near coords, point meta=explicit symbolic,
    error bars/.cd, y dir=both, y explicit,
  ] table[x=molecule, y=naive, y error minus=naive_em, y error plus=naive_ep,
          meta=naive_lbl] {figures/data/energy_methods_cluster.dat};
  \addlegendentry{naive GPU $O(NM)$}
  \addplot+[
    ybar, bar shift=14.0pt, fill=energyfmm, draw=energyfmm!70!black, line width=0.4pt,
    nodes near coords, point meta=explicit symbolic,
    error bars/.cd, y dir=both, y explicit,
  ] table[x=molecule, y=fmm, y error minus=fmm_em, y error plus=fmm_ep,
          meta=fmm_lbl] {figures/data/energy_methods_cluster.dat};
  \addlegendentry{FMM}

  \coordinate (dnf-H1N1-cpu) at ([xshift=-14.0pt]axis cs:H1N1,20000);

\end{axis}
\draw[-{Stealth[length=2.6mm,width=2.2mm]}, line width=1pt, energycpu!75!black]
  (dnf-H1N1-cpu) -- ++(0,0.62cm)
  node[above, inner sep=1pt, font=\scriptsize\bfseries, text=black!70]{DNF};

\end{tikzpicture}

\caption[Energy stage cost by method]{\textbf{Cost of the energy stage on
    \energyMoleculesenergymethodscluster{}.} Median wall time of the energy calculation stage on a
    logarithmic axis; the number above each bar is that median in
    seconds.\energyMachineenergymethodscluster\energyCaveatsenergymethodscluster}
\label{fig:energy-methods-cluster}
\end{figure}

%  Needs \usepackage{pgfplots} and \usetikzlibrary{arrows.meta} in the preamble.
% ------------------------------------------------------------------------
% generated by scripts/plot_energy_methods.py -- do not edit
\def\energyMoleculesenergymethodscluster{1VSZ, 6VYB, H1N1}
\def\energyRepeatsenergymethodscluster{3}
\def\energyMachineenergymethodscluster{ Measured on an A100 node of the TU Berlin HPC (Table~\ref{tab:test-systems}) with 16 MPI ranks.}
\def\energyFmmConfigenergymethodscluster{ The FMM ran at $P = 9$, $\theta = 0.3$, $n_{leaf,src} = 16$, $n_{leaf,tgt} = 1024$ for 1VSZ and 6VYB; $P = 10$, $\theta = 0.4$, $n_{leaf} = 256$ for H1N1.}
\def\energyReferenceenergymethodscluster{the CPU path where it ran and the naive GPU sum otherwise}
\def\energyCaveatsenergymethodscluster{ H1N1: measured on four H200 GPUs of the TU Berlin HPC (Table~\ref{tab:test-systems}) with 8 MPI ranks; 1 run per method, not 3.}


\pgfplotsset{compat=1.16}

\definecolor{energycpu}{HTML}{D95F02}
\definecolor{energynaive}{HTML}{9DC3E0}
\definecolor{energyfmm}{HTML}{2C5F8D}

\begin{figure}[!htb]
\centering
\begin{tikzpicture}
\begin{axis}[
    ybar,
    % 15cm against the 16.03cm textwidth of the 12pt/2.5cm page. Measured, not
    % assumed: 15.4 still fits, 13.6 gains nothing.
    width=15.0cm, height=7.4cm,
    bar width=13pt,
    enlarge x limits=0.30,
    % A shared linear axis would put every bar but the slowest on the baseline:
    % the times span 5756x. log origin y=infty is what makes the bars start at
    % ymin rather than at 1, which on a log axis is where ybar puts them by
    % default -- without it the sub-second bars would hang below the frame.
    ymode=log, log origin y=infty,
    ymin=0.05, ymax=20000,
    ytick={0.1,1,10,100,1000,10000},
    % 1, 10, 100 rather than 10^0, 10^1, 10^2. The times are wall clock and get
    % read off as such; the exponents make a reader translate before comparing.
    log ticks with fixed point,
    ylabel={energy stage time [s]},
    ylabel style={font=\footnotesize},
    symbolic x coords={1VSZ,6VYB,H1N1},
    xtick=data,
    % The atom count belongs on the axis, not in the caption: the three columns
    % are three problem sizes, and without it the H1N1 column reads as the same
    % work taking longer rather than as two orders of magnitude more of it.
    xticklabels={{1VSZ\\{\scriptsize 14\,299 atoms}},{6VYB\\{\scriptsize 44\,495 atoms}},{H1N1\\{\scriptsize 14\,442\,610 atoms}}},
    xticklabel style={font=\footnotesize, align=center},
    tick label style={font=\scriptsize},
    ymajorgrids=true,
    major grid style={black!12},
    % Clipping stays ON. log origin y=infty gives every bar a base far below the
    % frame -- a large finite log coordinate, not a real infinity -- and the clip
    % is the only thing keeping those bases off the page. The DNF annotation
    % therefore cannot be drawn in here; it is anchored to a \coordinate and
    % drawn after \end{axis}, where clipping no longer applies.
    % -0.22, not the -0.14 the other figures use: the tick labels here are two
    % lines, and the legend has to clear the atom count as well as the name.
    legend style={
      at={(0.5,-0.22)}, anchor=north, draw=none,
      legend columns=3, font=\footnotesize, column sep=0.8em,
    },
    error bars/error bar style={black!85, line width=0.8pt},
    error bars/error mark options={
      mark size=3pt, rotate=90, line width=0.8pt, black!85,
    },
    % The white backing is load-bearing on the H1N1 group: a four-digit label is
    % wider than a 13pt bar, so it overhangs onto the neighbouring bar, which
    % there is a full-height block of colour.
    every node near coord/.append style={
      font=\tiny, color=black!55, anchor=south, yshift=4pt,
      inner sep=1pt, fill=white, fill opacity=0.85, text opacity=1,
    },
]
  \addplot+[
    ybar, bar shift=-14.0pt, fill=energycpu, draw=energycpu!70!black, line width=0.4pt,
    nodes near coords, point meta=explicit symbolic,
    error bars/.cd, y dir=both, y explicit,
  ] table[x=molecule, y=cpu, y error minus=cpu_em, y error plus=cpu_ep,
          meta=cpu_lbl] {figures/data/energy_methods_cluster.dat};
  \addlegendentry{CPU}
  \addplot+[
    ybar, bar shift=0.0pt, fill=energynaive, draw=energynaive!70!black, line width=0.4pt,
    nodes near coords, point meta=explicit symbolic,
    error bars/.cd, y dir=both, y explicit,
  ] table[x=molecule, y=naive, y error minus=naive_em, y error plus=naive_ep,
          meta=naive_lbl] {figures/data/energy_methods_cluster.dat};
  \addlegendentry{naive GPU $O(NM)$}
  \addplot+[
    ybar, bar shift=14.0pt, fill=energyfmm, draw=energyfmm!70!black, line width=0.4pt,
    nodes near coords, point meta=explicit symbolic,
    error bars/.cd, y dir=both, y explicit,
  ] table[x=molecule, y=fmm, y error minus=fmm_em, y error plus=fmm_ep,
          meta=fmm_lbl] {figures/data/energy_methods_cluster.dat};
  \addlegendentry{FMM}

  \coordinate (dnf-H1N1-cpu) at ([xshift=-14.0pt]axis cs:H1N1,20000);

\end{axis}
\draw[-{Stealth[length=2.6mm,width=2.2mm]}, line width=1pt, energycpu!75!black]
  (dnf-H1N1-cpu) -- ++(0,0.62cm)
  node[above, inner sep=1pt, font=\scriptsize\bfseries, text=black!70]{DNF};

\end{tikzpicture}

\caption[Energy stage cost by method]{\textbf{Cost of the energy stage on
    \energyMoleculesenergymethodscluster{}.} Median wall time of the energy calculation stage on a
    logarithmic axis; the number above each bar is that median in
    seconds.\energyMachineenergymethodscluster\energyCaveatsenergymethodscluster}
\label{fig:energy-methods-cluster}
\end{figure}
